\section{Limitaciones y conclusiones}

\subsection{Limitaciones}

\begin{itemize}[leftmargin=1.2em,itemsep=0.15em]
  \item El horizonte es finito (30 días tras 1 día de warm-up). Los intervalos
  de confianza cuantifican la variabilidad entre réplicas, pero no eliminan el
  sesgo por condiciones iniciales ni por eventos raros de cola.
  \item Se usa drenaje al final del horizonte para no censurar las alertas
  admitidas. Sin drenaje, las métricas de tiempo podrían subestimarse.
  \item Se supone independencia entre llegadas y servicios, estacionariedad de
  las tasas y ausencia de abandono, reintentos o vacaciones de analistas.
  \item La prioridad es no expropiativa; una alerta crítica puede esperar a que
  termine un servicio preventivo en curso.
  \item La escala de la Weibull no es su media: $E[S_2]\approx10{,}8329$ min.
  \item Los resultados corresponden a los escenarios modelados; cambios en la
  mezcla de alertas o en la capacidad requieren nuevas corridas.
\end{itemize}

\subsection{Conclusiones}

El simulador reproduce el M/M/1 dentro de los intervalos de confianza y
satisface la Ley de Little, lo que valida el motor y la estimación. En el SOC
base, la pérdida de alertas es baja ($0{,}385$\,\%), pero las alertas críticas
también se bloquean y la prioridad no elimina la espera. El Escenario A muestra
que un aumento del $40$\,\% en alertas preventivas multiplica la pérdida por
$5{,}8$ y deteriora fuertemente la congestión, aunque la utilización teórica no
llegue al $100$\,\%. El Escenario B demuestra que ampliar el buffer no sustituye
capacidad de servicio: bajo sobrecarga, el cuarto analista reduce más la
pérdida y, sobre todo, los tiempos de espera. La recomendación técnica es
contratar el cuarto analista manteniendo $K=8$.
